\section{Program modules}
\progname{} has three user modules. The corresponding input jobs are listed below.
\begin{enumerate}
\item \textbf{diabat.} This module contains the following diabatization methods.
  \begin{itemize}
  \item \texttt{\$diabat-fphd}. Fragment particle--hole densities (FPHD).
  \item \texttt{\$diabat-gmh}. Multistate generalized Mulliken--Hush (GMH).
  \item \texttt{\$diabat-fcd}. Multistate fragment charge difference (FCD).
  \item \texttt{\$diabat-fed}. Fragment excitation difference (FED).
  \item \texttt{\$diabat-fedfcd}. Multistate FED--FCD.
  \end{itemize}
%FPHD is a distinctive feature of \progname{}. It permits diabatization of multichromophoric systems without prescribing the types of quasi-diabatic states in advance, provided that the states can be distinguished by their fragment particle and hole distributions.
It also includes the \texttt{\$diabat-utils-phasefix} utility for adjusting wave-function phases relative to selected reference states. See Chapter~\ref{chap-diabat}.

\item \textbf{postorb.}
This module provides postprocessing and analysis of atomic and molecular orbitals.
It supports the calculation of various one-electron integrals, the conversion among different orbital-file formats, and the generation of real-space grid data for orbitals, electron densities, and spin densities.
See Chapter~\ref{chap-postorb}.

\item \textbf{postwfn.}
This module provides postprocessing and analysis of electronic-state wave functions based on one-particle density matrices, transition density matrices, and difference density matrices.
It supports natural-orbital analyses, the calculation of various electronic and transition properties, and the generation of real-space grid data for various densities, such as electron, transition, and difference densities.
Both adiabatic and quasi-diabatic states can be analyzed.
See Chapter~\ref{chap-postwfn}.

\end{enumerate}
